% qibao95/qibao95.tex —— 七宝中学高一上周末作业 9.5（试卷版/答案版）
% 试卷版：xelatex qibao95.tex
% 答案版：xelatex "\def\isanswer{1}\input{qibao95.tex}"

\documentclass[a4paper,11pt]{article}
\usepackage{common}

\begin{document}

\examtitlest{2026--2027 学年度七宝中学高一上周末作业 9.5}{集合与常用逻辑用语}

\bigsec{一、填空题}

\begin{enumerate}
\setcounter{enumi}{2}

\item 已知集合 $A$ 满足若 $n\in A$ 且 $n\in\Z$，则 $\tfrac{1}{n}\in A$，小张同学迅速得出 3 个结论：（1）$0\notin A$；（2）集合 $A$ 不可能是单元集；（3）当 $n$ 取遍可以取的所有数时，集合元素的个数一定是偶数. 其中正确结论的序号为 \blankline[2em].

\ans{（1）（3）}

\ifanswer
\solbegin
\step{对于 (1)，若 $0\in A$，则 $0$ 不能作为分母，故 $0\notin A$.}
\stepm{(1)}{正确；}
\step{对于 (2)，当 $n=1$ 时，$\tfrac{1}{n}\in A$，所以 $A=\{1\}$，故 (2) 错误；}
\step{对于 (3)，由 (2) 的结论，当 $n=1$ 时，$\tfrac{1}{n}\in A$，即集合 $A$ 一定包含元素 $1$.}
\step{同理当 $n=-1$ 时，$\tfrac{1}{n}=-1\in A$，所以集合 $A$ 一定包含元素 $-1$.}
\step{当 $n$ 取其他整数时，则其倒数必在集合 $A$ 中.}
\step{所以当 $n$ 取遍可以取的所有数时，集合 $A$ 的元素一定为偶数.}
\step{故 (3) 正确.}
\step{故正确结论的序号为 (1)(3).}
\solend

\fi
\item 已知 $M=\{a-3,\,2a-1,\,a^2+1\}$，$N=\{-2,\,4a-3,\,3a-1\}$，若 $M=N$，则实数 $a$ 的值为 \blankline[2em].

\ans{$1$}

\ifanswer
\solbegin
\step{因为 $M=N$，所以 $a-3=-2$ 时，$a=1$，$M=\{-2,1,2\}$，$N=\{-2,1,2\}$，满足 $M=N$；}
\step{$2a-1=-2$ 时，$a=-\tfrac{1}{2}$，$M=\{-\tfrac{7}{2},-2,\tfrac{5}{4}\}$，$N=\{-2,-5,\tfrac{5}{2}\}$，不满足 $M=N$.}
\step{所以 $a=1$.}
\solend

\fi
\item 集合 $A=\{x\mid -12\leqslant x\leqslant 15\}$，集合 $B=\{x\mid m+1\leqslant x\leqslant 2m-1\}$，若 $A\cap B\neq\varnothing$，则实数 $m$ 的取值范围是 \blankline[2em].

\ans{$\{m\mid 2\leqslant m\leqslant 14\}$}

\ifanswer
\solbegin
\step{先求解 $A\cap B=\varnothing$ 时 $m$ 的取值范围.}
\stepm{①}{当集合 $B=\varnothing$ 时，则 $m+1>2m-1$，解得 $m<2$.}
\stepm{②}{当集合 $B\neq\varnothing$ 时，则 $\begin{cases}m+1\leqslant 2m-1 \\ 2m-1<-12\text{ 或 }\begin{cases}m+1\leqslant 2m-1 \\ m+1>15\end{cases}\end{cases}$，解得 $m>14$.}
\step{综上，$A\cap B=\varnothing$ 时，$m<2$ 或 $m>14$.}
\step{所以 $A\cap B\neq\varnothing$ 时，$m$ 的取值范围是 $\{m\mid 2\leqslant m\leqslant 14\}$.}
\solend

\fi
\item 设 $A=\{x\mid x^2-3x-4=0\}$，$B=\{x\mid ax-1=0\}$，若 $A\cap B=B$，则实数 $a$ 的值可以为 \blankline[2em].

\ans{$0,\,-1,\,\tfrac{1}{4}$}

\ifanswer
\solbegin
\step{$A=\{-1,4\}$，因为 $A\cap B=B$，所以 $B\subseteq A$.}
\stepm{①}{$a=0$ 时，$B=\varnothing$，满足 $B\subseteq A$.}
\stepm{②}{$a\neq 0$ 时，$B=\{\tfrac{1}{a}\}$，则 $\tfrac{1}{a}=-1$ 或 $\tfrac{1}{a}=4$，所以 $a=-1$ 或 $\tfrac{1}{4}$.}
\step{所以 $a$ 的值可以为 $0,\,-1,\,\tfrac{1}{4}$.}
\solend

\fi
\item 集合 $A=\{x\mid x\geqslant -1\}$，集合 $B=\{x\mid a-1\leqslant x\leqslant a+2\}$，若 $A\cap B=\varnothing$，则实数 $a$ 的取值范围是 \blankline[2em].

\ans{$(-\infty,-3)$}

\ifanswer
\solbegin
\step{因为 $a-1<a+2$，所以集合 $B\neq\varnothing$，因为 $A\cap B=\varnothing$，所以 $a+2<-1$，解得 $a<-3$.}
\step{所以实数 $a$ 的取值范围为 $(-\infty,-3)$.}
\solend

\fi
\item 已知集合 $A=\{1,3,a^2\}$，$B=\{1,a+2\}$，若 $A\cup B=A$，则 $a=$ \blankline[2em].

\ans{$2$}

\ifanswer
\solbegin
\step{由题意得 $a+2=3$ 或 $a+2=a^2$.}
\step{若 $a+2=3$，此时 $a=1\Rightarrow a^2=1$，集合 $A$ 的元素有重复，不符合题意；}
\step{若 $a+2=a^2$，解得 $a=2$ 或 $a=-1$.}
\step{显然 $a=2$ 时符合题意，而 $a=-1\Rightarrow a^2=1$ 同上，集合 $A$ 的元素有重复，不符合题意.}
\step{故 $a=2$.}
\solend

\fi
\item 已知全集 $U=\{x\in\mathbb{N}\mid x\leqslant 9\}$，$A\cap B=\{3\}$，$\overline{A}\cap B=\{4,6,8\}$，$A\cap\overline{B}=\{1,5\}$，则 $\overline{A}=$ \blankline[2em].

\ans{$\{0,2,4,6,7,8,9\}$}

\ifanswer
\solbegin
\step{由题意得 $U=\{x\in\mathbb{N}\mid x\leqslant 9\}=\{0,1,2,3,4,5,6,7,8,9\}$，$A\cap B=\{3\}$ 得 $A$、$B$ 中含有元素 $3$；}
\step{$\overline{A}\cap B=\{4,6,8\}$ 得 $B$ 中含有 $4$、$6$、$8$，集合 $A$ 中不含 $4$、$6$、$8$；}
\step{$A\cap\overline{B}=\{1,5\}$ 得 $A$ 含有 $1$、$5$，集合 $B$ 不含 $1$、$5$.}
\step{对于 $0$，假设集合 $A$ 中含有 $0$，由 $A\cap B=\{3\}$ 得 $B$ 不含 $0$，则 $A\cap\overline{B}$ 必含元素 $0$，与 $A\cap\overline{B}=\{1,5\}$ 矛盾.}
\step{同理可得集合 $A$ 中不含元素 $2$、$7$、$9$.}
\step{综上所述，$A=\{1,3,5\}$.}
\step{故 $\overline{A}=\{0,2,4,6,7,8,9\}$.}
\solend

\fi
\item 已知集合 $A=\{x\mid x<-1\text{ 或 }x>4\}$，$B=\{x\mid 2a\leqslant x\leqslant a+3\}$，若 $B\subseteq A$，则实数 $a$ 的取值范围是 \blankline[2em].

\ans{$\{a\mid a<-4\text{ 或 }a\geqslant 2\}$}

\ifanswer
\solbegin
\step{若 $B\subseteq A$，则 $B=\varnothing$ 或 $B\neq\varnothing$.}
\step{当 $B=\varnothing$ 时，则 $2a>a+3$，即 $a>3$.}
\step{当 $B\neq\varnothing$ 时，$a\leqslant 3$ 且 $a+3<-1$ 或 $2a>4$，解得 $a<-4$ 或 $3\geqslant a>2$.}
\step{综上，$a<-4$ 或 $a\geqslant 2$.}
\solend

\fi
\item 已知集合 $M=\{m\mid$ 关于 $x$ 的方程 $\tfrac{x^2-9}{2x-m}=1$ 有唯一解$\}$，用列举法表示 $M=$ \blankline[2em].

\ans{$\{-6,6,10\}$}

\ifanswer
\solbegin
\step{由 $\tfrac{x^2-9}{2x-m}=1$ 得 $x^2-2x+m-9=0$，则 $\Delta=4-4\times 1\times(m-9)=0$，解得 $m=10$.}
\step{当方程 $\tfrac{x^2-9}{2x-m}=1$ 可化为二元一次方程时，此时二元一次方程有一解，符合题意.}
\step{由于 $x^2-9=(x+3)(x-3)$，则令 $x-\tfrac{m}{2}=x+3$ 或 $x-\tfrac{m}{2}=x-3$，解得 $m=-6$ 或 $m=6$.}
\step{此时 $\tfrac{x^2-9}{2x-m}=1$ 为 $\tfrac{x-3}{2}=1$ 或 $\tfrac{x+3}{2}=1$，只有一解，符合题意.}
\step{综上所述，集合 $M=\{-6,6,10\}$.}
\solend

\fi
\item 已知非空数集 $S$ 满足：对任意给定的 $x$，$y\in S$ ($x$、$y$ 可以相同)，有 $x+y\in S$ 且 $x-y\in S$. 若集合 $S$ 中最小的正数为 6，则集合 $S=$ \blankline[2em].

\ans{$\{t\mid t=6n,\,n\in\Z\}$}

\ifanswer
\solbegin
\step{由对任意给定的 $x$、$y\in S$ ($x$、$y$ 可以相同)，有 $x+y\in S$ 且 $x-y\in S$，又 $6$ 是集合中的最小正整数，则 $6n,\,n\in\Z$ 也在集合 $S$ 里.}
\step{假设 $S$ 里有形如 $6n+r,\,r\in\Z,\,r\in\{1,2,3,4,5\}$，那么 $(6n+r)-6n=r\in S$，与 $6$ 是集合中的最小正整数矛盾.}
\step{故 $S=\{t\mid t=6n,\,n\in\Z\}$.}
\solend

\fi
\end{enumerate}

\bigsec{二、选择题}

\begin{enumerate}
\setcounter{enumi}{12}

\item 集合 $A=\{x\in\mathbb{N}\mid\tfrac{12}{x}\in\mathbb{N}\}$ 的子集的个数为（\quad）.

A. $64$\qquad B. $16$\qquad C. $6$\qquad D. $4$

\ans{A}

\ifanswer
\solbegin
\step{由题意得 $x$ 的值可以为 $1,2,3,4,6,12$，所以 $A=\{1,2,3,4,6,12\}$ 有 6 个元素，所以集合 $A$ 的子集有 $2^6=64$ 个.}
\step{故选 A.}
\solend

\fi

\item 平面直角坐标系中，除去 $(1,2)$、$(3,4)$ 两点的所有点构成的集合为（\quad）.

A. $\{(x,y)\mid x\neq 1\text{ 且 }x\neq 3\}$；

B. $\{(x,y)\mid x\neq 1\text{ 或 }y\neq 2\text{ 或 }x\neq 3\text{ 或 }y\neq 4\}$；

C. $\{(x,y)\mid (x-1)^2+(y-2)^2\neq 0\text{ 且 }(x-3)^2+(y-4)^2\neq 0\}$；

D. $\{(x,y)\mid x\neq 1,\,y\neq 2,\,x\neq 3,\,y\neq 4\}$.

\ans{C}

\item 已知 $m\in\R$，集合 $A=\{x\mid m<x<3m-1\}$ 中的元素恰有 2 个整数，则 $m$ 的取值范围是（\quad）.

A. $\{2\}$\qquad B. $\left[\tfrac{4}{3},\tfrac{5}{3}\right)$\qquad C. $\R$\qquad D. $\left[\tfrac{4}{3},\tfrac{5}{3}\right)\cup\{2\}$

\ans{D}

\ifanswer
\solbegin
\step{因为集合 $A=\{x\mid m<x<3m-1\}$ 中的元素恰有两个整数，所以 $1<3m-1-m\leqslant 3$，解得 $1<m\leqslant 2$.}
\stepm{①}{当 $1<m<2$ 时，集合 $A$ 中的两个整数分别为 $2$、$3$，则 $3<3m-1\leqslant 4$，解得 $\tfrac{4}{3}<m\leqslant\tfrac{5}{3}$.}
\stepm{②}{当 $m=2$ 时，此时 $A=\{x\mid 2<x<5\}$，符合题意.}
\step{综上所述，实数 $m$ 的取值范围是 $\left[\tfrac{4}{3},\tfrac{5}{3}\right)\cup\{2\}$.}
\step{故选 D.}
\solend

\fi

\item 已知集合 $M=\{a,b,c,d,e\}$，若集合 $N_1$、$N_2$、$\dots$、$N_k$ 是集合 $M$ 的 $k$ 个不同子集，且 $N_1\cup N_2\cup\dots\cup N_k\neq M$，则 $k$ 的最大值为（\quad）.

A. $15$\qquad B. $16$\qquad C. $31$\qquad D. $32$

\ans{B}

\ifanswer
\solbegin
\step{因为集合 $M=\{a,b,c,d,e\}$，共有 5 个元素，所以集合 $M$ 的子集个数为 $2^5=32$ 个.}
\step{因为 $N_1\cup N_2\cup\dots\cup N_k\neq M$，所以并集缺少 $M$ 中至少一个元素.}
\step{设缺少元素 $a$，则所有子集 $N_i$ 均不含 $a$，即 $N_i\subseteq\{b,c,d,e\}$.}
\step{集合 $\{b,c,d,e\}$ 的子集个数为 $2^4=16$，满足条件.}
\step{假设 $k\geqslant 17$，因为要满足这些子集的并集不等于 $M$，必须至少有一个 $M$ 中的元素不在任何一个子集中，否则并集就会等于 $M$.}
\step{设这个缺失的元素是 $a$，那么所有这 $k$ 个子集都不含 $a$，即每个子集都是 $\{b,c,d,e\}$ 的子集，而 $\{b,c,d,e\}$ 只有 $2^4=16$ 个不同的子集，我们却要从中取出至少 17 个不同的子集，这是不可能的.}
\step{因此假设不成立，所以 $k$ 不可能大于 16.}
\step{因此 $k$ 的最大值为 16.}
\step{故选 B.}
\solend

\fi
\end{enumerate}

\bigsec{三、解答题}

\begin{enumerate}
\setcounter{enumi}{16}

\item 已知集合 $A=\{x\mid x^2-3x+2=0\}$，集合 $B=\{x\mid x^2+2(a+1)x+a^2-5=0\}$.
\begin{enumerate}[label=(\arabic*)]
\item 若集合 $B$ 为单元集集合，求实数 $a$ 的值；
\item 若 $A\cup B=A$，求实数 $a$ 的取值范围.
\end{enumerate}

\ans{(1) $a=-3$；(2) $(-\infty,-3)$}

\ifanswer
\solbegin
\stepm{(1)}{因为集合 $B=\{x\mid x^2+2(a+1)x+a^2-5=0\}$ 仅有一个元素，所以 $\Delta=4(a+1)^2-4(a^2-5)=8a+24=0$，解得 $a=-3$.}
\stepm{(2)}{因为 $A\cup B=A$，所以 $B\subseteq A$.}
\stepm{①}{当 $B=\varnothing$ 时，则 $\Delta=8a+24<0$，解得 $a<-3$，符合题意.}
\stepm{②}{当 $B$ 有一个元素时，$a=-3$，此时 $x^2-4x+4=0$，所以 $B=\{2\}$，符合题意.}
\stepm{③}{当 $B=\{1,2\}$ 时，则 $\begin{cases}1+2=-2(a+1) \\ 1\times 2=a^2-5\end{cases}$，无解.}
\step{综上所述，$a$ 的取值范围为 $(-\infty,-3]$.}
\solend

\fi
\ansspace[6]
\item 已知集合 $A=\{x\mid 2<x<5\}$，$B=\{x\mid m-2\leqslant x\leqslant 2m-1\}$.
\begin{enumerate}[label=(\arabic*)]
\item 若 $m=5$，求 $B$；
\item 若 $B=\varnothing$，求实数 $m$ 的取值集合；
\item 若 $B$ 中有 3 个整数，求实数 $m$ 的取值集合；
\item 若 $A\subseteq B$，求实数 $m$ 的取值集合；
\item 若 $B\subseteq A$，求实数 $m$ 的取值集合.
\end{enumerate}

\ans{(1) $B=\{x\mid 3\leqslant x\leqslant 9\}$；(2) $\{m\mid m<-1\}$；(3) $[1,2)\cup\left(2,\tfrac{5}{2}\right]$；(4) $\{m\mid 3\leqslant m\leqslant 4\}$；(5) $\{m\mid m<-1\}$}

\ansspace[12]
\item 已知数集 $X=\{x_1,x_2,\dots,x_n\}$（其中 $1\leqslant x_1<x_2<\dots<x_n,\,n\in\mathbb{N}^*,n\geqslant 2$）具有性质 $M$：对任意的 $i,j\,(1\leqslant i\leqslant j\leqslant n)$，$x_ix_j$ 与 $\tfrac{x_i}{x_j}$ 两数中至少有一个属于 $X$.
\begin{enumerate}[label=(\arabic*)]
\item 分别判断数集 $P=\{1,2,3\}$ 与 $Q=\{1,2,4\}$ 是否具有性质 $M$；
\item 证明：$x_1=1$；
\item 已知数集 $X$ 具有性质 $M$，若 $n=5$，$x_2=3$，求数集 $X$.
\end{enumerate}

\ans{(1) $P$ 不具有性质 $M$，$Q$ 具有性质 $M$；(2) 略；(3) $X=\{1,3,9,27,81\}$}

\ifanswer
\solbegin
\stepm{(1)}{对于集合 $P$：取 $x_2=2,\,x_3=3$，则 $x_2x_3=2\times 3=6\notin P$，$\tfrac{x_3}{x_2}=\tfrac{3}{2}\notin P$，所以数集 $P$ 不具有性质 $M$.}
\step{对于集合 $Q$：$1\times 2=2\in Q$，$1\times 4=4\in Q$，$\tfrac{4}{2}=2\in Q$，$\tfrac{4}{4}=1\in Q$，则对 $Q$ 中任意 $x_i$、$x_j$，$x_ix_j$ 与 $\tfrac{x_i}{x_j}$ 中至少有一个属于 $Q$，所以数集 $Q$ 具有性质 $M$.}
\stepm{(2)}{由已知 $x_2\geqslant 1$.}
\step{若 $x_1>1$，则 $x_n>1$，$x_nx_n>x_n$，所以 $x_nx_n\notin X$，又 $\tfrac{x_n}{x_n}=1\notin X$，所以数集 $X$ 不具有性质 $M$，不符合题意.}
\step{所以 $x_1=1$.}
\stepm{(3)}{数集 $X$ 具有性质 $M$，$n=5$，$x_2=3$.}
\step{由 (2) 得 $x_1=1$.}
\step{因为 $1<x_2<x_3<x_4<x_5$，所以 $x_i\cdot x_5>x_i\,(i=2,3,4,5)$，所以 $x_i\cdot x_5\notin X$.}
\step{因为数集 $X$ 具有性质 $M$，所以 $\tfrac{x_5}{x_1}\in X$，所以 $x_5=1=\tfrac{x_5}{x_5}<\tfrac{x_5}{x_4}<\tfrac{x_5}{x_3}<\tfrac{x_5}{x_2}<x_5$，所以 $\tfrac{x_5}{x_4}=x_2$，$\tfrac{x_5}{x_3}=x_3$，$\tfrac{x_5}{x_2}=x_4$，所以 $x_5=x_2x_4=x_3x_4$，即 $x_5^2=x_2x_4^2$.}
\step{因为 $x_2x_4>x_3x_4=x_5$，所以 $x_3x_4\notin X$.}
\step{所以 $\tfrac{x_4}{x_3}\in X$，由 $1<\tfrac{x_2}{x_3}<\tfrac{x_2}{x_2}=x_3<1<\tfrac{x_2}{x_2}<x_3<x_5$，所以 $\tfrac{x_4}{x_3}=x_2$，所以 $\tfrac{x_5}{x_4}=\tfrac{x_5}{x_3}=\tfrac{x_5}{x_2}=x_5$，所以 $x_3=x_2^2=9,\,x_4=x_2x_3=3\times 9=27,\,x_5=x_3^2=9^2=81$.}
\step{所以 $X=\{1,3,9,27,81\}$.}
\solend

\fi
\ansspace[12]
\ansneed{7cm}
\item 已知集合 $M=\{a_1,a_2,a_3,\dots,a_n\}(0\leqslant a_1<a_2<a_3<\dots<a_n,\,n\in\mathbb{N}^*)$ 具有性质：对任意 $1\leqslant i\leqslant j\leqslant n\,(i,j\in\mathbb{N}^*)$，$a_i+a_j$ 与 $a_j-a_i$ 至少有一个属于 $M$，则称 $M$ 为“封闭集”.
\begin{enumerate}[label=(\arabic*)]
\item 若集合 $P=\{2,5,8\}$，$Q=\{0,5,10\}$，判断 $P$、$Q$ 是否是“封闭集”，并说明理由；
\item 若集合 $P=\{a_1,a_2,a_3\}(0\leqslant a_1<a_2<a_3)$ 是“封闭集”，且 $a_3=2026$，求集合 $P$；
\item 设集合 $A=\{a_1,a_2,a_3,\dots,a_n\}(0\leqslant a_1<a_2<a_3<\dots<a_n,\,n\geqslant 4,\,n\in\mathbb{N}^*)$ 是“封闭集”，证明：$a_1+a_2+a_3+\dots+a_n=\tfrac{n}{2}\cdot a_n$.
\end{enumerate}

\ans{(1) $P$ 不是“封闭集”，$Q$ 是“封闭集”；(2) $\{0,1013,2026\}$；(3) 略}

\ifanswer
\solbegin
\stepm{(1)}{集合 $P$ 不是“封闭集”，理由如下：在集合 $P=\{2,5,8\}$ 中，因为 $2+5=7\notin P$，$5-2=3\notin P$，所以集合 $P$ 不是“封闭集”.}
\step{集合 $Q$ 是“封闭集”，理由如下：集合 $Q=\{0,5,10\}$ 中，因为 $0+5=5\in Q$，$0+10=10\in Q$，$10-5=5\in Q$，$0+0=0\in Q$.}
\step{$5+5=10\in Q$，$5-5=0\in Q$，$10-10=0\in Q$.}
\step{所以集合 $Q$ 是“封闭集”.}
\stepm{(2)}{因为 $a_1<a_2<a_3$，且 $P=\{a_1,a_2,a_3\}$ 是“封闭集”，$a_3+a_3>a_3$，所以 $a_3+a_3\notin P$，则 $a_3-a_3=0\in P$，所以 $a_1=0$.}
\step{因为 $a_2+a_3>a_3$，所以 $a_2+a_3\notin P$.}
\step{则 $(a_3-a_3)\in P,(a_3-a_2)\in P,(a_2-a_2)\in P,(a_3-a_1)\in P,(a_2-a_1)\in P$.}
\step{$a_3-a_3<a_3-a_2<a_2-a_1$.}
\step{由集合的元素互异性得 $a_3-a_3=a_2$，因为 $a_3=2026$，所以 $a_2=\tfrac{1}{2}a_3=1013$.}
\step{所以集合 $P=\{0,1013,2026\}$.}
\stepm{(3)}{因为 $M=\{a_1,a_2,a_3,\dots,a_n\}(0\leqslant a_1<a_2<a_3<\dots<a_n,\,n\in\mathbb{N}^*)$ 是“封闭集”，所以 $a_n+a_n\notin A$，则 $a_n-a_n=0\in A$，所以 $a_1=0$.}
\step{又 $0\leqslant a_1<a_2<\dots<a_n$，所以 $0\leqslant a_n-a_n<a_n-a_{n-1}<\dots<a_n-a_1$.}
\step{又 $a_n+a_{n-1}>a_n\,(i=1,2,\dots,n-1)$，所以 $a_n+a_{n-1}\notin A$，则 $a_n-a_{n-1}\in A$.}
\step{由集合元素的互异性得 $a_1=a_n-a_n,\,a_2=a_n-a_{n-1},\,a_3=a_n-a_{n-2},\,\dots,\,a_n=a_n-a_1$.}
\step{所以 $a_1+a_2+a_3+\dots+a_n=(a_n-a_1)+(a_n-a_{n-1})+(a_n-a_{n-2})+\dots+(a_n-a_1)=na_n-(a_1+a_2+a_3+\dots+a_n)$.}
\step{所以 $2(a_1+a_2+a_3+\dots+a_n)=na_n$，所以 $a_1+a_2+a_3+\dots+a_n=\tfrac{n}{2}\cdot a_n$.}
\solend

\fi
\ansspace*[11]
\end{enumerate}

\end{document}